\chapter{无限集}

\section{自然数集}

\begin{definition}{无限集,无限基数}{}
	不是有限集的集合称为无限集.不是有限基数的基数称为无限基数.
\end{definition}

\begin{axiom}{无穷公理}{}
	存在无限集.
\end{axiom}

\begin{theorem}{自然数集的存在性}{}
	公式``$x$是自然数''是集化的.
\end{theorem}

\begin{definition}{自然数集及其基数}{}
	我们把所有自然数组成的集合记作$\N$,称之为自然数集,把$\left|\N\right|$记作$\aleph_{0}$.没有强调时,$\N$上的偏序按照基数的偏序理解.
\end{definition}

\begin{definition}{序列,无限序列,集合$E$中的序列}{}
	以$\N$的子集为指标集的族称为序列,以$\N$的无限子集为指标集的序列称为无限序列,集合$E$中以$\N$的子集为指标集的族称为集合$E$中的序列.
\end{definition}

\begin{convention}{}{}
	设$P\left\{i\right\}$是公式,满足$P\left\{i\right\}$的元素$i$构成一个由自然数组成的有限集.我们通常把$E$中的序列$\left(x_{i}\right)_{i\in I}$写作$\left(x_{i}\right)_{P\left\{i\right\}}$.特别地,如果$I=\left[a,b\right]$,则把$\left(x_{i}\right)_{i\in I}$记作$\left(x_{i}\right)_{a\leq i\leq b}$或$\left(x_{i}\right)_{i=a}^{b}$.
	
	同样的,对于集合族$\left(X_{i}\right)_{i\in I}$的乘积,相应地使用$\prod\limits_{P\left\{i\right\}}X_{i}$或$\prod\limits_{n=k}^{\infty}X_{i}$.类似的记号对$\left(X_{i}\right)_{i\in I}$的并,交,基数和,基数积都适用.
\end{convention}

\begin{definition}{子序列}{}
	序列的子族称为给定序列的子序列.
\end{definition}

\begin{definition}{只有项的顺序不同的序列}{}
	设$\left(x_{i}\right)_{i\in I},\left(y_{i}\right)_{i\in I}$是集合$E$中的序列.若存在置换$f:I\to I$使得$\forall n\in I\left(x_{f\left(n\right)}=y_{n}\right)$,则称$\left(x_{i}\right)_{i\in I}$和$\left(y_{i}\right)_{i\in I}$只有项的顺序不同.
\end{definition}

\begin{definition}{$p$重序列}{}
	设$p\in\N$.以$\underbracket{\N\times\ldots\times\N}_{p\text{个}}$为指标集的序列称为$p$重序列.
\end{definition}

\begin{definition}{族的序列化}{}
	设集合$I$与$\N$等势,$f:\N\to I$为双射.对任意序列$\left(x_{i}\right)_{i\in I}$,称序列$\left(x_{f\left(n\right)}\right)_{n=0}^{\infty}$为通过$f$诱导的序定义的序列.
\end{definition}

\begin{remark}{}{}
	对给定的自然数$n$,通过上述方法可以从指标集是$n$元集的族得到一个指标集为$\left[1,n\right]$或$\left[0,n-1\right]$的序列.
\end{remark}
























